% ========================================================================
% فصل ۵: مکانیزم‌های حفظ حریم خصوصی
% خروجی عامل sam-master: sam-work/sections/05-privacy.tex
% منابع مجاز: rahmawatiagustinaSystematicLiteratureReview2025،
%   luPseudonymChangingSocial2012، SecuringVehicularAd،
%   farsimadanReviewSecurityChallenges2025، VanetSecurityPrivacy،
%   zhangRAISEEfficientRSUAided2008، mazzocca2025survey
% ========================================================================

\section{مکانیزم‌های حفظ حریم خصوصی}

پس از بررسی حملات امنیتی و به‌ویژه حملات بر حریم خصوصی در فصل پیشین، این فصل به مکانیزم‌های مقابله با آن‌ها می‌پردازد. شبه‌نام و چرخهٔ حیات آن، راهبردهای تغییر شبه‌نام و تکنیک‌های بی‌نامی، سه رکن اصلی این مکانیزم‌ها هستند که هر یک در یک زیربخش بررسی می‌شود.

\subsection{شبه‌نام و چرخه حیات آن}

شبه‌نام\LTRfootnote{\lr{Pseudonym}} شناسه‌ای موقت است که هویت واقعی خودرو را پنهان می‌کند و به خودرو امکان مشارکت در شبکه را می‌دهد \cite{rahmawatiagustinaSystematicLiteratureReview2025}. در همین مرور، شبه‌نام‌ها شناسه‌هایی موقت هستند که به خودروها اجازه می‌دهند به‌صورت ناشناس ارتباط برقرار کنند و در عین حال تعاملات امن و احرازهویت‌شده باقی بمانند \cite{rahmawatiagustinaSystematicLiteratureReview2025}. ویژگی دوم شبه‌نام، پیوندناپذیری\LTRfootnote{\lr{Unlinkability}} آن است: چون خودرو در طول مسیر از شبه‌نام‌های مختلفی استفاده می‌کند، پیوندناپذیری این شبه‌نام‌ها حریم خصوصی موقعیت را تضمین می‌کند \cite{luPseudonymChangingSocial2012}. در این تعریف باید دقت شود که ویژگی «موقت» بودن به مرور راهماتی و همکاران و ویژگی پیوندناپذیری به لو و همکاران نسبت دارد؛ هیچ‌یک از این دو منبع هر دو صفت را در یک تعریف واحد جمع نکرده است.

انگیزهٔ این مکانیزم آشکار است: تبادل مداوم داده در شبکه به مهاجم امکان می‌دهد از خودروها نمایه‌سازی\LTRfootnote{\lr{Profiling}} کند و بدین‌سان حریم خصوصی مالکان آن‌ها را به خطر بیندازد \cite{rahmawatiagustinaSystematicLiteratureReview2025}. از سوی دیگر، حریم خصوصی هویت پیش‌شرط حریم خصوصی موقعیت است و به همین دلیل هر خودرو باید برای پخش پیام‌ها به جای هویت واقعی از شبه‌نام استفاده کند \cite{luPseudonymChangingSocial2012}. تعادل میان گمنامی و پاسخ‌گویی\LTRfootnote{\lr{Accountability}} نیز در همین مرور بازتاب می‌یابد: خودروها بدون افشای هویت واقعی خود به نهادهای غیرمجاز تأیید می‌شوند، اما نهادهای مجاز می‌توانند در صورت نیاز خودروهای مخرب را ردیابی و شناسایی کنند \cite{rahmawatiagustinaSystematicLiteratureReview2025}.

مدیریت شبه‌نام فرایندی پنج‌مرحله‌ای است که در چرخهٔ حیات شبه‌نام\LTRfootnote{\lr{Pseudonym life cycle}} از صدور تا لغو را در بر می‌گیرد \cite{rahmawatiagustinaSystematicLiteratureReview2025}. این چرخه شامل مراحل زیر است:

\begin{enumerate}
  \item \textbf{صدور\LTRfootnote{\lr{Issuance}}:} به خودروها شبه‌نام‌های امن از نظر رمزنگاری تخصیص می‌یابد.
  \item \textbf{استفاده\LTRfootnote{\lr{Usage}}:} ارتباطی ناشناس اما تأییدشده از طریق شبه‌نام ممکن می‌شود.
  \item \textbf{تغییر\LTRfootnote{\lr{Changing}}:} شبه‌نام‌ها به‌طور منظم تغییر می‌کنند تا از ردیابی پرهیز شود.
  \item \textbf{حل\LTRfootnote{\lr{Resolution}} (بازگشایی هویت):} نهادهای مجاز در سناریوهای حادثهٔ امنیتی هویت واقعی پشت یک شبه‌نام را شناسایی می‌کنند.
  \item \textbf{لغو\LTRfootnote{\lr{Revocation}}:} خودروهای آسیب‌دیده از اتصال به شبکه منع می‌شوند تا ریسک و دسترسی ناخواسته کاهش یابد.
\end{enumerate}

در مرحلهٔ صدور، تولید شبه‌نام یا توسط یک نهاد انجام می‌شود یا با همکاری چند نهاد؛ در حالت همکاری، رایج‌ترین ترکیب‌ها خودرو و مرجع قابل اعتماد\LTRfootnote{\lr{Trusted Authority (TA)}} با ۴۶ درصد و خودرو و مرجع ردیابی\LTRfootnote{\lr{Trace Authority (TRA)}} با ۳۵ درصد هستند \cite{rahmawatiagustinaSystematicLiteratureReview2025}. توابع درهم‌ساز\LTRfootnote{\lr{Hash function}} به دلیل سربار محاسباتی کم، رایج‌ترین روش تولید شبه‌نام‌اند و در حدود ۷۳ درصد مطالعات بررسی‌شده به کار رفته‌اند؛ روش‌های دیگر رمزنگاری و امضای دیجیتال هستند \cite{rahmawatiagustinaSystematicLiteratureReview2025}. برخی طرح‌ها از مُهرِ زمانی اعتبار\LTRfootnote{\lr{Validity timestamp}} استفاده می‌کنند تا تغییر شبه‌نام روان‌تر انجام شود \cite{rahmawatiagustinaSystematicLiteratureReview2025}.

در مرحلهٔ استفاده، شبه‌نام در ارتباطات خودرو به خودرو (\lr{V2V}) و خودرو به زیرساخت (\lr{V2I}) برای اشتراک اطلاعات ترافیکی، ارسال هشدارهای ایمنی و انجام وظایف رانندگی همکارانه به کار می‌رود \cite{rahmawatiagustinaSystematicLiteratureReview2025}. در مرحلهٔ حل، نهادهای مجاز مانند نیروهای انتظامی با استفاده از اطلاعات رمزنگاری‌شده‌ای که نزد مرجع قابل اعتماد یا مرجع ردیابی نگهداری می‌شود، هویت واقعی پشت یک شبه‌نام را در موارد حادثهٔ امنیتی شناسایی می‌کنند؛ در این مرحله مرجع قابل اعتماد یا مرجع ردیابی به‌عنوان واسط عمل می‌کند تا اطلاعات فقط به نهادهای مجاز افشا شود \cite{rahmawatiagustinaSystematicLiteratureReview2025}. در مرحلهٔ لغو نیز نهادی مجاز مانند مرجع قابل اعتماد، بر اساس شواهد تأییدشده از رفتار نادرست خودرو، اعلان لغو صادر می‌کند \cite{rahmawatiagustinaSystematicLiteratureReview2025}.

با وجود اهمیت این چرخه، پوشش مراحل در مطالعات بررسی‌شده نابرابر است. همهٔ ۱۱۳ مطالعهٔ اولیه از شبه‌نام استفاده کرده‌اند، اما مرحلهٔ تغییر تنها در ۲۱ درصد (۲۴ مورد)، مرحلهٔ حل در ۸۳ درصد (۹۴ مورد) و مرحلهٔ لغو در ۴۳ درصد (۴۹ مورد) پیاده شده است \cite{rahmawatiagustinaSystematicLiteratureReview2025}. در مجموع، تنها ۱۴ مورد از ۱۱۳ مطالعه، یعنی حدود ۱۲ درصد، هر پنج مرحله را پیاده‌سازی کرده بودند \cite{rahmawatiagustinaSystematicLiteratureReview2025}. پیامد این ناقص بودن جدی است: بیشتر طرح‌ها فقط بر صدور و استفاده تمرکز کرده‌اند و بدون پیاده‌سازی کامل چرخه، خودروها در برابر حملهٔ پیوند\LTRfootnote{\lr{Linkability attack}}، حملهٔ سیبل\LTRfootnote{\lr{Sybil attack}} و دسترسی غیرمجاز آسیب‌پذیر می‌مانند \cite{rahmawatiagustinaSystematicLiteratureReview2025}. نبود مرحلهٔ حل مسئلهٔ مرجع حل هویت\LTRfootnote{\lr{Privacy-authority problem}} را ایجاد می‌کند و نبود مرحلهٔ لغو خطر حملهٔ سیبل و دسترسی غیرمجاز را بالا می‌برد \cite{rahmawatiagustinaSystematicLiteratureReview2025}.

\subsection{استراتژی‌های تغییر شبه‌نام \lr{(ETSI TR 103 415)}}

پوشش کم مرحلهٔ تغییر در مطالعات بررسی‌شده، اهمیت انتخاب راهبرد درست برای این مرحله را برجسته می‌کند. مرحلهٔ تغییر با هدف جلوگیری از ردیابی بلندمدت انجام می‌شود و گزارش فنی \lr{ETSI TR 103 415} شش راهبرد برای آن برشمرده است \cite{rahmawatiagustinaSystematicLiteratureReview2025}. در مقدمهٔ همین مرور هشدار داده شده که راهبرد ثابت‌زمانی که چگالی ترافیک را نادیده بگیرد، به حملهٔ پیوند می‌انجامد \cite{rahmawatiagustinaSystematicLiteratureReview2025}. این شش راهبرد عبارت‌اند از:

\begin{itemize}
  \item \textbf{پارامتر ثابت\LTRfootnote{\lr{Fixed parameter}}:} شبه‌نام بر اساس یک پارامتر ثابت تغییر می‌کند.
  \item \textbf{تصادفی\LTRfootnote{\lr{Randomness}}:} شبه‌نام با فواصل زمانی تصادفی تغییر می‌کند.
  \item \textbf{دورهٔ سکوت\LTRfootnote{\lr{Silent period}}:} ارسال پیام در بازهٔ تغییر شبه‌نام به‌طور موقت متوقف می‌شود.
  \item \textbf{خودرومحور\LTRfootnote{\lr{Vehicle-centric}}:} تصمیم به تغییر بر اساس وضعیت خودرو گرفته می‌شود.
  \item \textbf{مبتنی بر چگالی\LTRfootnote{\lr{Density-based}}:} شبه‌نام بر اساس چگالی خودروهای اطراف تغییر می‌کند.
  \item \textbf{ناحیهٔ اختلاط\LTRfootnote{\lr{Mix zone}}:} شبه‌نام همهٔ خودروها در یک منطقه به‌صورت هماهنگ تغییر می‌کند تا پیوند پیام‌های متوالی قطع شود \cite{VanetSecurityPrivacy}.
\end{itemize}

پراکندگی این راهبردها در ۲۴ مطالعه‌ای که مرحلهٔ تغییر را پیاده کرده‌اند نیز نابرابر است: هشت مورد راهبرد خودرومحور، هفت مورد پارامتر ثابت، یک مورد تصادفی، دو مورد ترکیب خودرومحور و تصادفی، یک مورد ترکیب خودرومحور و پارامتر ثابت و پنج مورد سایر راهبردها را به کار برده‌اند \cite{rahmawatiagustinaSystematicLiteratureReview2025}. در جدول ۱۱ این مرور هیچ مقاله‌ای ذیل دورهٔ سکوت، مبتنی بر چگالی یا ناحیهٔ اختلاط فهرست نشده و اینکه آیا این سه راهبرد ذیل «سایر موارد» قرار گرفته‌اند مشخص نیست \cite{rahmawatiagustinaSystematicLiteratureReview2025}.

یک چالش متقاطع در همهٔ این راهبردها، پیش‌بینی‌پذیری\LTRfootnote{\lr{Predictability}} شبه‌نام است. در مرور راهماتی و همکاران هشدار داده شده که ورودی‌های ثابت تولید شبه‌نام، یعنی هویت واقعی و مُهر زمانی و عدد تصادفی، پیوند دادن شبه‌نام‌ها را آسان می‌کند، به‌ویژه اگر بازهٔ تغییر شبه‌نام قابل پیش‌بینی باشد \cite{rahmawatiagustinaSystematicLiteratureReview2025}. این مرور پیشنهاد می‌کند که ورودی‌های پویایی مانند شرایط جاده، مسیر خودرو و موقعیت جغرافیایی افزوده شود، هرچند باید میان پیش‌بینی‌ناپذیری و ردیابی‌پذیری\LTRfootnote{\lr{Traceability}} تعادل برقرار شود \cite{rahmawatiagustinaSystematicLiteratureReview2025}.

نمونه‌ای از پیاده‌سازی این راهبردها، تغییر شبه‌نام در نقاط اجتماعی\LTRfootnote{\lr{Social spot}} است که لو و همکاران پیشنهاد می‌کنند. نقاط اجتماعی جایگاه‌هایی هستند که چندین خودرو در آن‌ها به‌طور موقت جمع می‌شوند، مانند تقاطع هنگام چراغ قرمز یا پارکینگ رایگان نزدیک یک مرکز خرید \cite{luPseudonymChangingSocial2012}. اگر همهٔ خودروها در این نقاط به‌طور نامتمایز شبه‌نام‌های خود را تغییر دهند، نقاط اجتماعی به‌طور طبیعی به ناحیهٔ اختلاط تبدیل می‌شوند \cite{luPseudonymChangingSocial2012}. در این طرح، نخستین پیام پس از تغییر برای همهٔ خودروها یکسان است: موقعیت برابر نقطهٔ اجتماعی، سرعت برابر صفر و شبه‌نامی پیوندناپذیر \cite{luPseudonymChangingSocial2012}.

\subsection{تکنیک‌های بی‌نامی}

راهبردهای تغییر شبه‌نام تنها بخشی از راهکارها هستند و با تکنیک‌های بی‌نامی\LTRfootnote{\lr{Anonymity}} تقویت می‌شوند که پیوند پیام‌ها به یک خودروی خاص را از ریشه دشوار می‌کنند. در مرور فرسیمادان و همکاران، تحلیل ترافیک\LTRfootnote{\lr{Traffic analysis}} به‌عنوان حمله‌ای منفعل در دستهٔ حملات خودخواهانه بررسی شده و از جمله دفاع‌ها برای آن کلیدهای ناشناس پویا\LTRfootnote{\lr{Dynamic anonymous keys}}، شبه‌نام و امضای گروهی\LTRfootnote{\lr{Group signature}} برشمرده شده است \cite{farsimadanReviewSecurityChallenges2025}.

رایا و هوبو به جای اصطلاح شبه‌نام از «جفت‌کلید ناشناس\LTRfootnote{\lr{Anonymous key pair}}» بهره می‌گیرند: جفت‌کلیدی که مرجع صدور گواهی\LTRfootnote{\lr{Certificate Authority (CA)}} آن را تأیید کرده، اما نه اطلاعاتی دربارهٔ هویت واقعی خودرو (یعنی پلاک الکترونیکی\LTRfootnote{\lr{Electronic License Plate (ELP)}}) دارد و نه ارتباط آشکاری با آن؛ با این حال این ناشناسی برای تعیین مسئولیت، مشروط است \cite{SecuringVehicularAd}. خودرو مجموعه‌ای از این کلیدها را نگه می‌دارد تا ردیابی نشود \cite{SecuringVehicularAd}، زیرا با ثبت پیام‌های یک کلید می‌توان فرستنده را تا کشف هویتش ردیابی کرد، مثلاً از روی محل سکونت \cite{SecuringVehicularAd}. این ناشناسی مشروط توزیع‌شده نیز هست: توان افشای هویت باید میان چند مرجع تقسیم شود، به‌طوری که پلیس بدون اجازهٔ قاضی نتواند هویت را بازیابی کند، و رایا و هوبو برای دسترسی به پایگاه‌دادهٔ ارتباط پلاک الکترونیکی و کلیدهای ناشناس، تسهیم راز\LTRfootnote{\lr{Secret sharing}} \lr{Shamir} را پیشنهاد می‌کنند \cite{SecuringVehicularAd}.

مرور سیلا و همکاران ناحیه‌های اختلاط را منطقه‌هایی می‌داند که در آن خودروها به‌صورت هماهنگ شبه‌نام‌های خود را تغییر می‌دهند تا پیوند پیام‌های متوالی قطع شود \cite{VanetSecurityPrivacy}. این مرور تغییر شبه‌نام در مکان‌ها یا بازه‌های زمانی راهبردی، احراز هویت حافظ حریم خصوصی\LTRfootnote{\lr{Privacy-preserving authentication}} و اثبات دانش‌صفر\LTRfootnote{\lr{Zero-knowledge proof}} (پراوین، ۲۰۲۱) را از جمله راهکارهای حفظ گمنامی و حریم خصوصی می‌شمارد \cite{VanetSecurityPrivacy}.

رویکرد دیگر، \lr{k}-گمنامی\LTRfootnote{\lr{k-anonymity}} است. در طرح \lr{RAISE}، که احراز اصالت پیام با کمک واحد کنارجاده‌ای\LTRfootnote{\lr{RSU}} است، نویسندگان از رویکرد \lr{k}-گمنامی برای حفاظت از حریم خصوصی هویت کاربر استفاده می‌کنند، به‌طوری که یک مهاجم نمی‌تواند یک پیام را با یک خودروی خاص مرتبط کند \cite{zhangRAISEEfficientRSUAided2008}. در این طرح، واحدهای کنارجاده‌ای مسئول تأیید اصالت پیام‌های ارسالی از خودروها و اطلاع‌رسانی نتیجه به خودروها هستند \cite{zhangRAISEEfficientRSUAided2008}.

اثبات دانش‌صفر در مرور مازوکا و همکاران یکی از دسته‌های افشای انتخابی\LTRfootnote{\lr{Selective disclosure}} برای اعتبارنامه‌های تأییدپذیر\LTRfootnote{\lr{Verifiable Credential (VC)}} است؛ راه‌حل پیشروی کنونی، اعتبارنامهٔ \lr{SD-JWT} است که با جایگزینی ادعاهای متن‌صریح با خلاصه‌های مالح‌دار\LTRfootnote{\lr{Salted digest}} آن‌ها، حریم خصوصی را تقویت می‌کند \cite{mazzocca2025survey}. در بخش حمل‌ونقل هوشمند، این مرور به کاربرد اعتبارنامه‌های مبتنی بر اثبات دانش‌صفر در ارتباطات خودرویی غیرمتمرکز (\lr{D-V2X}) اشاره می‌کند \cite{mazzocca2025survey}.
